% Report-ready mathematical section. Requires amsmath, amssymb, amsthm.
% All source-specific input names and figures should be supplied by the main report.
\section{Why the two logical changes compose}
\label{sec:composition-audit}

There are two different mathematical comparisons in this construction.
Changing the output-root K schema to the seed-request schema preserves
the relevant consistency strength.  Merging the equality row with its
helper preserves the entire native theorem set.  The second comparison
does not depend on which set-theoretic schema is present, so it applies
to the seed-request theory without repeating or strengthening its
set-theoretic interpretation.

\subsection{The exact theories being compared}

Write $I(A,B)$ for the native implication constructor, $Q_i A$ for the
native universal constructor, and $d_{ab}$ for its primitive equality
atom.  The designated refutation target is the literal native code
$F=1$.  Put $N(A)=I(A,F)$ and $X_iA=N(Q_iN(A))$.  These definitions are
operations on natural-number codes, including the exceptional small
codes; they are not an assumption that the total decoder is a globally
injective constructor homomorphism.

Let
\[
 \kappa(a,b)=4\left(\frac{(a+b)(a+b+1)}2+b\right)+1.
\]
The seed predecessor contains the old equality row and helper
\[
 \rho_7(a,b,c)=d_{ca}\to(d_{cb}\to d_{ab}),
 \qquad E(a,b)=d_{ab}\to Q_{\kappa(a,b)}d_{ab}.
\]
Its actual enabled bundle has the grouping
\[
 B=N\bigl(I(V,I(E,T))\bigr),\qquad
 T=Q_c I(H,N(W)),
\]
where $V$ is membership vacuity and $H,W$ are the exact inherited seed
closure and witness request.  All four expressions depend on the actual
source parameters.  The new source replaces $\rho_7$ by
\[
 G(a,b,c)=d_{ca}\to
       \bigl(d_{cb}\to Q_{\kappa(a,b)}d_{ab}\bigr)
\]
and replaces the enabled bundle by
\[
 B'=N\bigl(I(V,T)\bigr).
\]
It retains the same gate $c\ne a,b$, all the other logical rows, the
truth seed, modus ponens and theorem generalization.  The row $G$ is
ungated and ranges over all triples, as did the deleted $\rho_7$.

\subsection{A closed native derivation of the missing equality facts}

The unchanged $\rho_8$ supplies the classical implicational basis.
The unchanged $\rho_6(i,A)$ is $X_iQ_iA\to A$.  Because $X_iQ_iA$
has the form $U\to F$, the implicational basis gives
$F\to X_iQ_iA$, and hence $F\to A$.  Consequently ordinary Boolean
tautologies with $F$ false can be used as abbreviations for finite
native derivations.  This argument makes no freshness assumption about
the label $0$ occurring in the decoded atom for $F$.

From a closed theorem $A\to B$, theorem generalization followed by
$\rho_2$ gives $Q_iA\to Q_iB$.  Classical contraposition gives
$X_iA\to X_iB$.  Here a closed theorem means a derivation without open
assumptions; its formula may have free numerical variable labels.

The following two implications use only the retained rows:
\[
 Q_iA\to X_iA, \qquad Q_iQ_iA\to A.                 \tag{C1}
\]
For the first, use $\rho_1(i,i)$, namely $X_id_{ii}$.
The quantified version of $A\to(d_{ii}\to A)$ and the normal
quantifier-distribution rule give
$Q_iA\to Q_i(d_{ii}\to A)$, while the existential consequence of
that distribution gives
$Q_i(d_{ii}\to A)\to(X_id_{ii}\to X_iA)$.
The second implication follows by applying the first to $Q_iA$ and
then $\rho_6(i,A)$.  This does not assume the unrestricted principle
$Q_iA\to A$.

Fix $x$, let $z=\kappa(x,x)$, and set $e=d_{xx}$.
The instance $G(x,x,z)$, followed by propositional contraction, gives
\[
 d_{zx}\to Q_ze.
\]
Existential monotonicity and $\rho_1(z,x)$ yield $X_zQ_ze$;
$\rho_6(z,e)$ then yields $d_{xx}$.  Thus primitive reflexivity is
derived without using either of the deleted axioms.  Next,
$G(a,b,a)$ and this reflexivity give the exact old helper
\[
                 d_{ab}\to Q_{\kappa(a,b)}d_{ab}.       \tag{C2}
\]
Generalize (C2) at $i=\kappa(a,b)$ and apply $\rho_2$:
\[
 Q_i d_{ab}\to Q_iQ_i d_{ab}.
\]
Composing with (C1) gives $Q_i d_{ab}\to d_{ab}$.  Composing this
last implication with $G(a,b,c)$ recovers the literal old
$\rho_7(a,b,c)$.  In the reverse direction, the old $\rho_7$ and
the old helper immediately give $G$ by propositional composition.

Every generalization above is applied to a theorem.  No quantifier is
introduced while retaining an open hypothesis, no equality-substitution
rule is silently added, and no unrestricted primitive equality-vacuity
schema is assumed.

\subsection{The seed bundle has its own exact propositional certificate}

The three needed implications are
\[
 B\to B',\qquad B\to E,\qquad B'\to(E\to B).           \tag{C3}
\]
They are Boolean tautologies for the displayed native grouping.
For clarity, $V,E,T$ can be treated as three arbitrary propositional
atoms: $B$ means $V\land E\land\neg T$, while $B'$ means
$V\land\neg T$.  Their exact code expressions are retained in (C3);
no equality of differently associated numerical constructor terms is
claimed.

Every pair $(a,b)$ obtains its old helper from an enabled old bundle:
choose $c=\max(a,b)+1$ and apply the second implication in (C3).
The helper binder is $\kappa(a,b)$, independent of that gate parameter.
Thus every new $G$ row is derivable in the old seed theory.  Conversely,
the new $G$ rows supply (C2) for every pair, and the third implication in
(C3) recovers each enabled old bundle from its corresponding new bundle.
Disabled triples continue to emit the same truth seed.

Write $\operatorname{Thm}_{\mathrm{s},E}(n)$ and
$\operatorname{Thm}_{\mathrm{s},G}(n)$ for derivability of the native
code $n$ in these two theories.  Replace each changed axiom leaf by the
finite closed derivation above, preserve MP and GEN, and induct on the
finite proof tree.  This establishes
\[
 \boxed{\quad
 \operatorname{Thm}_{\mathrm{s},E}(n)
 \ \longleftrightarrow\
 \operatorname{Thm}_{\mathrm{s},G}(n)
 \quad\text{for every natural }n.\quad}
\]
The substitutions are effective syntactic operations, the gate witness
is explicit, and the proof-tree induction is an elementary
arithmetical argument formalizable in ZFC.  In particular the
distinguished target $n=1$ is preserved literally.  Composing this
equivalence with the inherited seed-theory calibration gives the same
internal ZFC-consistency equivalence for the new numerical source.

There is a useful generalization for future source searches.  The same
interderivation works for any total label function $\kappa(a,b)$:
its use of $z=\kappa(x,x)$ never assumes that $z\ne x$.  To retain the
inherited standard-equality consistency calibration, it suffices that
$a\ne b$ imply $\kappa(a,b)\notin\{a,b\}$; there is no restriction
on diagonal values.  This is precisely the general-helper hypothesis
proved in the retained Source116 adapter.  All eight native
\texttt{Imp}/\texttt{Rel}/\texttt{All}/\texttt{Eq} builders with
either ordering of their two arguments satisfy it.  In particular,
the implication builder's sole collision at $(0,0)$ is permitted.
The released construction uses the displayed membership-code binder;
this general logical statement is not an additional state-count claim
for any of the other builders.

\subsection{What carries over and what must be checked again}

The inherited seed calibration treats $c$ simultaneously as matrix code
and private container label.  Its soundness proof uses the exact total
decoder, the special cases $D(0)$ and $D(1)$, and the theorem
$c\notin\operatorname{FV}(D(c))$.  At $c=0$ the required statement
is closedness; it is incorrect to assert that every bound name in
$D(0)$ is strictly below zero.  Distinct active labels use the ordinary
two-role abstraction; coincident active labels use the separate
empty-set/singleton argument.  The lower interpretation translates
equality into membership indiscernibility before choosing the padded
matrix's native code as its private name.  The present literal
theorem-set comparison leaves all these arguments unchanged.

The conclusion about the binary machine additionally requires that the
new emitter really produces $B'$ and $G$ with the stated parameter
coverage, that its source enumeration is correct, that source-tree
changes preserve execution, and that the final transition table
simulates this source from the exact blank initialization.  The source
and table checks in the main report discharge those separate
obligations.  Neither propositional certificates nor bounded simulation
alone discharge them.

The combined finite proof DAG has 318 nodes, 293 formulas and 36 roots.
It consists of the predecessor's 315-node G proof DAG and the three
additional seed-order tautologies (C3).  An independent parser checks
each MP edge, each theorem-GEN node, each permitted axiom instance and
each truth-table tautology.  Its forward G roots contain neither an
old-helper nor an old-$\rho_7$ leaf; its reverse roots contain no G leaf.
These finite instances accompany the uniform proof just given.  They
are not a newly compiled end-to-end Rocq or Lean theorem for the final
machine bytes.
